% ======================================================================
% 第 II 节：夸克束缚机制
% ======================================================================
\section{夸克束缚机制}\label{sec:binding}

本节将用费曼路径积分图像说明束缚机制。路径积分框架将按直觉而非严格的方式使用。
考虑流-流传播子
\begin{equation}
D_{\mu\nu}(x) = \ev{\mathrm{T}\,J_{\mu}(x)J_{\nu}(0)}\,,
\label{eq:2.1}
\end{equation}
其傅里叶变换决定 $e^{+}e^{-}$ 湮灭为强子的截面。假设流 $J_{\mu}(x)$ 由夸克场
构成，如同夸克-部分子模型那样。假设夸克通过单一规范场相互作用。（限制为一个
场只是为了简便。）在费曼路径积分图像中，传播子 $D_{\mu\nu}(x)$ 是对所有可能
的经典夸克路径以及规范场的所有可能经典值取加权平均给出的。流 $J_{\mu}(x)$ 和
$J_{\nu}(0)$ 被认为在原点产生一个夸克对，之后在点 $x$ 处湮灭：必须对这一对
夸克中的每一个都对连接 $0$ 和 $x$ 两点的所有路径求和。

图~\ref{fig:3} 展示了夸克与反夸克路径的一个例子。真空也能发射和吸收夸克对；
这导致图~\ref{fig:4} 所画的额外闭圈。所有可能的圈也必须求和。对于点 $0$ 和
$x$ 也可能存在独立的圈（图~\ref{fig:5}），但这种可能性在此并不重要。

\begin{figure}[t]
  \centering
  \includegraphics[width=0.52\textwidth]{fig3}
  \refstepcounter{figure}%
  \caption*{图 3. 连接 $0$ 和 $x$ 两点的夸克 ($q$) 与反夸克 ($\bar q$) 路径
    的一个例子。}
  \label{fig:3}
\end{figure}

\begin{figure}[t]
  \centering
  \includegraphics[width=0.5\textwidth]{fig4}
  \refstepcounter{figure}%
  \caption*{图 4. 带额外真空圈（如图 3 中那样）的流圈的例子。}
  \label{fig:4}
\end{figure}

与给定的夸克路径或路径组相联系的权重包含因子
$\exp[ig\oint A_{\mu}(x)\,ds^{\mu}]$，其中 $A_{\mu}(x)$ 是规范场。这里
$\oint ds^{\mu}$ 是对每个夸克-反夸克圈（包括连接 $0$ 与 $x$ 的圈）的线积分或
线积分之和。常数 $g$ 是规范理论的耦合常数。还有与 $A$ 无关的其它权重因子。
最后，与夸克路径无关，还有另一个权重因子，即规范场自由作用的指数。然后将
组合权重因子对所有夸克路径和所有规范场 $A_{\mu}(x)$ 求平均，得到流-流传播子。

\begin{figure}[t]
  \centering
  \includegraphics[width=0.42\textwidth]{fig5}
  \refstepcounter{figure}%
  \caption*{图 5. 点 $0$ 与 $x$ 处独立夸克圈的例子。（对规范场积分会产生连接
    这些圈的规范传播子。）}
  \label{fig:5}
\end{figure}

为了使夸克作为独立的末态粒子存在，必须可能出现夸克与反夸克线分离得很开的
夸克-反夸克圈，至少在 $x$ 与 $0$ 相距很远时要如此。图~\ref{fig:6}(a) 说明了
这一点。如果夸克与反夸克路径不太可能分离到超过一个固定的大小，比如
$10^{-13}$~cm [见图~\ref{fig:6}(b)]，那么显然不会有探测器看到孤立的夸克或
反夸克。

在本讨论中假定真空圈不重要。如果真空圈足够重要，那么空间将被高密度的真空产生
的夸克-反夸克对所填满。特别是，在宏观大小的探测器内部会有许多夸克-反夸克对。
问题随之变成：在一个宏观大小的区域内能否出现夸克对反夸克的过程（或反之）。
这是一个更难回答的问题，本文不作讨论。

\begin{figure}[t]
  \centering
  \includegraphics[width=0.55\textwidth]{fig6}
  \refstepcounter{figure}%
  \caption*{图 6. (a) 夸克与反夸克分离良好的圈。(b) 夸克与反夸克间距很小的圈。}
  \label{fig:6}
\end{figure}

注意 $x$ 必须很大：如果 $x$ 很小，找到大尺度圈的可能性很小。这看起来有些奇怪：
人们预期，如果夸克在 $e^{+}e^{-}$ 湮灭末态中出现，那也只会在虚光子动量 $q$ 很大
时出现，而大的 $q$ 意味着小的 $x$，而不是大的 $x$。这一悖论的答案有两部分。
首先，费曼路径积分中的重要路径与可能的物理末态没有细致关系（这些路径是裸粒子的
路径，而非物理粒子的路径）。其次，大的 $x$ 不一定意味着小的 $q$。事实上，研究
大 $x$ 时是否出现分离良好的夸克-反夸克路径，实际上是在寻找 $e^{+}e^{-}$ 湮灭中的
夸克-反夸克阈，它将给流-流传播子贡献一项 $\sim \exp(2imx)/x^{p}$（对大 $x$），
其中 $m$ 是夸克质量。（这里 $\sim$ 表示至多相差 $x$ 的幂次。${}^{9}$）这样一项
在动量空间对应阈处的奇点
$q^{2} = (2m)^{2}$。

假设先对规范场求平均，再进行夸克路径平均。那么需要确定权重因子
$\exp[ig\oint A_{\mu}(y)\,ds^{\mu}]$ 再乘上自由规范场作用的指数之后对一切
规范场的平均。对于阿贝尔规范理论，这一平均可以显式计算：它等于
\[
  \exp\!\Big[-g^{2}\oint\!\!\oint ds^{\mu}\,ds'^{\nu}\,D_{\mu\nu}(y-y')\Big],
\]
其中 $D_{\mu\nu}(y-y')$ 是自由规范场传播子。

将上述表达式在一空间维与三空间维情形下作比较，可以看到夸克束缚机制。在三空间维
中，这一计算不给出束缚（束缚只在修改的规范场作用下出现：见第四节），而在
一空间维中存在束缚。在三空间维中，$D_{\mu\nu}(y-y')$ 的行为如同 $(y-y')^{-2}$。
因此 $(y'-y)$ 的大值在二重线积分中可忽略。于是二重积分正比于 $P$，其中 $P$ 是
圈长。不幸的是积分在 $y'=y$ 处发散；为使积分有意义需要截断。由于本文反正要引入
截断，这并非大问题。对于简单圈，周长 $P$ 大致为 $(x^{2})^{1/2}$ 量级（不考虑
$x$ 靠近光锥的情形）。于是得到人们预期自由夸克存在时的指数型行为。

在一空间维中，对于大的 $y'-y$，$D_{\mu\nu}(y-y')$ 的行为如同
$\ln[(y-y')^{2}]$，并且 $y'$ 与 $y$ 可以在圈上自由地独立取值。在这种情况下
二重积分正比于 $P^{2}$。于是规范场平均的行为如同 $e^{-cP^{2}}$，其中 $c$ 是常数。
在这种情况下大圈的贡献被强烈压低，不存在自由夸克。[图~\ref{fig:6}(b) 中夸克与
反夸克紧邻的情形是特殊的——此时大的 $y-y'$ 不重要，因为夸克路径与邻近的反向
反夸克路径之间相互抵消。这种情况下二重积分的行为如同 $P$ 而非 $P^{2}$，但此时
不存在孤立夸克。]

在本文后面描述的强耦合格点规范理论中，$\exp[ig\oint A_{\mu}(x)\,ds^{\mu}]$
的规范场平均的行为如同 $\exp(-cA)$，其中 $A$ 是圈所围的面积。这强烈压低大的圈，
例如图~\ref{fig:6}(a) 中的圈，其中 $A$ 是 $P^{2}$ 量级。可以认为一个因子 $P$
大致是 $(x^{2})^{1/2}$，另一个 $P$ 类似于一个质量乘以 $(x^{2})^{1/2}$。由于
$P \to \infty$ 当 $x \to \infty$，夸克-反夸克阈位于无穷大质量处。

\begin{figure}[t]
  \centering
  \includegraphics[width=0.42\textwidth]{fig7}
  \refstepcounter{figure}%
  \caption*{图 7. 邻近存在真空圈的夸克-反夸克圈。}
  \label{fig:7}
\end{figure}

在所有上述计算中，如果附近存在真空圈（图~\ref{fig:7}），则可以有大的圈。
此时总能得到 $e^{-cP}$ 行为。例如，在强耦合情形下，相关的围面积是两圈之间的
面积，只要两圈的距离独立于 $P$ 固定，该面积就正比于周长 $P$。这与施温格的图像
一致。孤立的、分离良好的圈可能被压低（由于指数中的 $P^{2}$ 或 $A$ 依赖），而
被真空极化圈紧密屏蔽的圈则从不被压低。

这里提出的束缚机制是软的：指数阻尼与大尺度、大面积的圈相联系。小圈的行为与
束缚机制无关。而且对小圈而言，其面积和周长都很小，在指数中都不太重要。

这里讨论的机制对狄拉克夸克或标量夸克同样适用。这与希格斯机制形成对比，后者
只能消除标量粒子的电荷。
